\documentclass{article}
\usepackage[UTF8]{ctex}
\usepackage{amsmath,amssymb,amsfonts}
\usepackage{enumitem}
\usepackage{geometry}
\geometry{a4paper, margin=1in}
% === 错题本排版优化（不参与题目标识） ===
\setlength{\parindent}{0pt}
\setlength{\parskip}{0.35em plus 0.10em minus 0.08em}
\setlength{\emergencystretch}{3em}
\setlist{itemsep=0.25em, topsep=0.40em, parsep=0pt}
\allowdisplaybreaks[2]
% === 错题本排版优化结束 ===
\begin{document}
\section*{华中科技大学2024年数学分析真题}
\section*{132．华中科技大学2024年数学分析试卷}

\subsection{题目1}
1．计算极限 \(\lim _{x \rightarrow 0^{+}} \frac{1}{x} \arccos \frac{\sin x}{x}\) ．

\subsection{题目2}
2．计算极限 \(\lim _{n \rightarrow \infty}\left(1+\frac{1}{n}\right)\left(1+\frac{1}{2^{3} n}\right)^{2^{3}}\left(1+\frac{1}{3^{3} n}\right)^{3^{3}} \cdots\left(1+\frac{1}{n^{3} n}\right)^{n^{3}}\)

\begin{itemize}

\subsection{题目3}
\end{itemize}

\subsection{题目3}
原题边界“3．”：已知 \(f\) 在 \(\mathbb{R}\) 上具有连续导数，\(L\) 是从 \(A\left(3, \frac{2}{3}\right)\) 到 \(B(1,2)\) 的曲线段 \(y=-\frac{5}{9} x^{2}+\frac{14}{9} x+1\) ，求
\[
I=\int_{L} \frac{1+y^{2} f(x y)}{y} \mathrm{~d} x+\frac{x}{y^{2}}\left(y^{2} f(x y)-1\right) \mathrm{d} y .
\]

\subsection{题目4}

\subsection{题目4}
原题边界“4．”：计算极限 \(\lim _{n \rightarrow \infty} \frac{1}{n} \int_{0}^{n \pi} \frac{x}{1+n^{2} \cos ^{2} x} \mathrm{~d} x\) ．

\subsection{题目5}
5．用下列两种方法求 \(I=\iint_{S} h(x, y, z) \mathrm{d} S\) ，其中 \(S: x^{2}+y^{2}+z^{2}=1, h(x, y, z)=x^{4}+y^{2} z^{2}\) 。

\begin{itemize}
  \item[（1）] 用球坐标变换．
  \item[（2）] 用 Gauss 公式和齐次函数 Euler 公式
\[
x f_{x}+y f_{y}+z f_{z}=n f
\]
其中 \(f\) 满足 \(f(t x, t y, t z)=t^{n} f(x, y, z), n\) 为正整数。
\end{itemize}

\subsection{题目6}
6．解答如下问题：

\begin{itemize}
  \item[（1）] 设 \(f\) 为 \([1,+\infty)\) 上的非负递减函数，记 \(x_{n}=\sum_{k=1}^{n} f(k)-\int_{1}^{n} f(t) \mathrm{d} t\) ，证明 \(\left\{x_{n}\right\}\) 收敛．
  \item[（2）] 计算 \(\sum_{n=1}^{\infty}(-1)^{n-1} \frac{\ln n}{n}\)（可能会用到 \(\sum_{k=1}^{n} \frac{1}{k}=\gamma+\ln n+o(1)(n \rightarrow \infty)\) ，其中 \(\gamma\) 为常数）．
\end{itemize}

\subsection{题目7}
7．设正整数 \(n, k\) 满足 \(n \geq k\) ，记 \(b(n, k)=\int_{\lambda}^{1} C_{n}^{k} t^{k-1}(1-t)^{n-k} \mathrm{~d} t\) ，其中 \(C_{n}^{k}=\frac{n!}{k!(n-k)!}\) ．

\begin{itemize}
  \item[（1）] 证明： \(\lim _{n \rightarrow \infty} \sum_{k=1}^{n} b(n, k)=2 \ln 2\) ．
  \item[（2）] 求极限 \(\lim _{n \rightarrow \infty} \frac{n 4^{n}}{3^{n}}\left(2 \ln 2-\sum_{k=1}^{n} b(n, k)\right)\) ．
\end{itemize}

\subsection{题目8}
8．设 \(u_{n}(x), S(x)\) 均为 \([a, b]\) 上的可积函数，其中 \(n=1,2, \cdots\) ．

\begin{itemize}
  \item[（1）] 证明：对 \(S_{n}(x)=\sum_{k=1}^{n} u_{k}(x)\) 有
\[
\lim _{n \rightarrow \infty} \int_{a}^{b}\left|S_{n}(x)-S(x)\right| \mathrm{d} x=0
\]
则逐项积分得到的级数 \(\sum_{n=1}^{\infty} \int_{a}^{x} u_{n}(t) \mathrm{d} t\) 在 \([a, b]\) 上一致收敛于 \(\int_{a}^{x} S(t) \mathrm{d} t\) 。
  \item[（2）] 说明函数项级数 \(\sum_{n=0}^{\infty} x^{n} \sin (\pi x)\) 在 \([0,1]\) 上不一致收敛，用（1）证明
\[
\sum_{n=0}^{\infty} \int_{0}^{1} x^{n} \sin (\pi x) \mathrm{d} x=\int_{0}^{1} \frac{\sin (\pi x)}{x} \mathrm{~d} x .
\]
\end{itemize}

\subsection{题目9}
9．设 \(a_{0}=1, a_{n+1}=a_{n} e^{\frac{1}{a_{n}}}, p>2\) ，证明：级数 \(\sum_{n=2}^{\infty} \frac{a_{n}-n}{n(\ln n)^{p}}\) 收敛。

\subsection{题目10}
10．设不恒为 0 的函数 \(f(x)\) 在 \([0,1]\) 上连续，且 \(\int_{0}^{1} f(x) \mathrm{d} x=0\) ．证明：存在 \(c \in(0,1)\) ，使得

\[
c^{2} f(c)=\int_{0}^{c}\left(\alpha x+x^{2}\right) f(x) \mathrm{d} x, 0<\alpha<\frac{4}{3} .
\]

\subsection{题目11}
11．设常数 \(\beta \geq 1\) ，证明对任意的 \(x>0\) ，成立

\[
\frac{x^{\beta}}{\sqrt{3}} \geq \frac{1}{x+\sqrt{3}-1}+\frac{1+\sqrt{3} \beta}{3} \ln x .
\]

\subsection{题目12}
12．解答如下问题：\\
（1）证明二元函数 \(z=f(P)\) 在 \(D \subset \mathbb{R}^{2}\) 上一致连续的充分必要条件是对 \(D\) 上任意两点列 \(\left\{P_{n}\right\}\) ，\(\left\{Q_{n}\right\}\) ，只要距离 \(d\left(P_{n}, Q_{n}\right) \rightarrow 0(n \rightarrow \infty)\) ，就有 \(f\left(P_{n}\right)-f\left(Q_{n}\right) \rightarrow 0(n \rightarrow \infty)\) ．\\
（2）判断函数

\[
g(x, y)= \begin{cases}\frac{x^{4} y^{4}}{x^{4}+y^{4}}, & x^{4}+y^{4} \neq 0 \\ 0, & x^{4}+y^{4}=0\end{cases}
\]

在 \(\mathbb{R}^{2}\) 上的一致连续性．
\end{document}
